% Generated draft: requires mathematical and source review.
\documentclass[11pt]{article}
\usepackage[margin=25mm]{geometry}
\usepackage{fontspec}
\setmainfont{Latin Modern Roman}
\usepackage{amsmath,amssymb,mathtools,unicode-math}
\setmathfont{Latin Modern Math}
\usepackage{enumitem,needspace,xurl,hyperref}
\setlength{\parindent}{0pt}
\setlength{\parskip}{5pt}
\setlength{\emergencystretch}{3em}
\newcommand{\N}{\mathbb{N}}
\newcommand{\Z}{\mathbb{Z}}
\newcommand{\Q}{\mathbb{Q}}
\newcommand{\R}{\mathbb{R}}
\newcommand{\C}{\mathbb{C}}
\newcommand{\F}{\mathbb{F}}
\newcommand{\A}{\mathbb{A}}
\newcommand{\PP}{\mathbb{P}}
\newcommand{\E}{\mathbb{E}}
\newcommand{\eps}{\varepsilon}
\newcommand{\dd}{\,\mathrm{d}}
\begin{document}
\section*{UnsolvedMath ID: 30004390}
\textbf{Source problem code:} catalogue entry
\subsection{Definitions and mathematical statement}
Let \(X\) and \(Y\) be spaces in the stated geometric setting: finite simply connected piecewise Riemannian simplicial complexes, or compact Riemannian manifolds. The supplied record does not specify whether simple connectivity is also required in the manifold case, or whether additional restrictions distinguish the source \(X\) from the target \(Y\).

Write \([X,Y]\) for homotopy classes of continuous maps \(X\to Y\). The record does not specify whether homotopies are based or unbased. For maps \(f,g\colon X\to Y\), a homotopy from \(f\) to \(g\) is a continuous map \(H\colon X\times[0,1]\to Y\) satisfying
\[
H(x,0)=f(x),\qquad H(x,1)=g(x)\quad(x\in X).
\]

The quantitative questions can be expressed schematically as follows. Once a geometric complexity functional \(G\), an algebraic size \(A\) on homotopy classes, and a homotopy complexity functional \(K\) have been specified, put
\[
\mathcal C(\alpha)=\inf_{[f]=\alpha}G(f),\qquad \alpha\in[X,Y],
\]
and, for homotopic maps \(f,g\), put
\[
\mathcal H(f,g)=\inf_{H\colon f\simeq g}K(H).
\]
These expressions are templates for the quantities requested; the supplied source does not define \(G\), \(A\), or \(K\), their admissible map classes, or any metric convention on \(X\times[0,1]\). Using infima does not assert that optimal representatives exist. The questions ask how to estimate \(\mathcal C(\alpha)\) in terms of \(A(\alpha)\), and how to estimate \(\mathcal H(f,g)\) in terms of \(G(f)\) and \(G(g)\). The intended bounds, asymptotic regimes, and permitted dependence on \(X\) and \(Y\) are unspecified.

Two further questions in the supplied formulation are:
\begin{enumerate}
\item Is scalability invariant under real homotopy equivalence? That is, for admissible spaces \(Y\) and \(Z\) of the same real homotopy type, must \(Y\) be scalable if and only if \(Z\) is scalable?
\item For an admissible space \(Y\), is the following equivalence valid?
\[
Y\text{ is scalable}
\quad\Longleftrightarrow\quad
\exists N\in\Z_{\geq0}\ \exists\text{ an injective differential-graded-algebra morphism }
H^*(Y;\R)\longrightarrow\bigwedge^*\R^N.
\]
Here \(H^*(Y;\R)\) has its cohomological grading, cup product, and zero differential; \(\bigwedge^*\R^N\) is the exterior algebra with generators in degree one and zero differential. Such a morphism preserves degrees, products, and differentials.
\end{enumerate}
The source does not define scalability or specify the intended notion of real homotopy type. Consequently, these two questions cannot be made fully self-contained from the supplied evidence. The exterior-algebra question occurs in the supplied statement and clean statement but is absent from the supplied original statement; its provenance requires review.

The supplied dated curation labels the problem partially solved and reports that scalability is a real-homotopy invariant and that equivalent algebraic criteria have been established. These are supplied literature assessments, not independently verified conclusions of this draft; in particular, the record does not explicitly identify the exact algebraic criterion proved.
\subsection{Short English statement}
The problem asks how the algebraic size of a homotopy class controls the least geometric complexity of a representative, and how the complexities of two homotopic maps control the least complexity of a homotopy between them. It also asks whether scalability depends only on real homotopy type and whether scalability is characterized by an embedding of real cohomology into a finite exterior algebra. The source leaves the complexity measures and scalability undefined. Its supplied dated assessment labels the problem partially solved; this draft does not independently verify that assessment.
\subsection{Sources}
\begin{itemize}
\item Supplied local UnsolvedMath registry record, permanent ID \(30004390\), title: \emph{Quantitative Complexity of Homotopy Representatives and Homotopies}; collection: Oberwolfach Reports Open Problems; exact source problem code: \texttt{catalogue entry}. No individual proposer is supplied. The record gives proposed year \(2020\).
\item Supplied source citation: \emph{Manifolds and Groups (2021)}, Oberwolfach Reports, pp.\ 417--463, DOI \texttt{10.4171/owr/2020/7}. \url{https://doi.org/10.4171/owr/2020/7}. The supplied background describes Oberwolfach Reports as documenting mathematical workshops and their open-problem discussions. The citation's year \(2021\), the proposed year \(2020\), and the DOI's year component are retained as supplied.
\item Supplied literature source: Berdnikov--Manin, \emph{Degrees of maps}. \url{https://doi.org/10.1017/fmp.2023.33}.
\item Literature assessment checked at \texttt{2026-08-22}: supplied dated curation, not independently verified. It labels the problem partially solved, describes the broad quantitative-homotopy questions as remaining active, and reports that scalability was proved to be a real-homotopy invariant and that equivalent algebraic criteria were established. The supplied background calls this dated literature triage derived from the OWR curation and requests independent verification of time-sensitive status and cited claims. Both supplied literature-source links are retained above.
\end{itemize}
\end{document}
